% Template de Artigo - Disciplina de Dados Categóricos
% Autora da disciplina: Profa. Ana Julia Alves Câmara (UFES)

\documentclass[12pt]{article}

% Pacotes básicos
\usepackage[utf8]{inputenc}
\usepackage[T1]{fontenc}
\usepackage[brazilian]{babel}
\usepackage{amsmath, amssymb}
\usepackage{graphicx}
\usepackage{booktabs}
\usepackage{setspace}
\usepackage{geometry}
\usepackage{hyperref}


\usepackage{tabularx}
\usepackage{makecell}

\geometry{margin=2.5cm}
\onehalfspacing

% Informações do artigo
\title{Título do Artigo}
\author{Nathália Nunes \and João Gabriel \and Rafael Jordane}
\date{Disciplina: Dados Categóricos \\
Departamento de Estatística -- UFES \\
\today}

\begin{document}

\maketitle

\begin{abstract}
Breve resumo do trabalho (150--200 palavras). O resumo deve apresentar:

\begin{itemize}
\item o problema de pesquisa
\item a base de dados utilizada
\item os métodos estatísticos aplicados
\item os principais resultados
\end{itemize}

\end{abstract}

\section{Pergunta}

A educação superior no Brasil cresce a cada ano, segundo dados do Censo Brasileiro da Educação Superior (2024) entre 2014 e 2024, a matrícula na educação superior aumentou 30,5\%, sendo 2,7\% o crescimento do ano de 2023 para 2024 de. Se em [1900\&bolinha] a formação se concentrava exclusivamente em áreas como [sla boticaria] [fonte], hoje cerca de [2000 mil] cursos de graduação diferentes são ofertados por instituições publicas e privadas no brasil [fonte]. O Censo da educação superior possui um importante papel de documentação desse crescimento, registrando com detalhes o perfil da educação brasileira a cada ano e possibilitando que esses dados sirvam como base para estudos acerca do tema, como foi feito nos artigos [A], [B] e [C]

Embora alguns autores como Mota (2012) tenham direcionado pesquisas em torno da projeção das oferta de vagas de cursos de áreas específicas como o Turismo, o campo social traz colaborações ainda mais abrangentes quando a disparidade da educação superior pública e privada são postas sob análise. Um exemplo disso é o tema da mercantilização da educação, tratada por Carvalho (2013), que ainda na década passada discutia o impacto do mercado financeiro e das grandes empresas na educação brasileira, visualizando que havia um contraponto entre o ensino público e privado.

Compreendendo complexidade da educação superior enquanto uma instituição social indivisível das demais questões que ceeca a sociedade brasileira, o presente estudo busca investigar se é possivel intentificar uma distinção clara entre as áreas dos cursos de graduação ofertados pelas instituições privadas e publicas no Brasil. O objetivo é que os resultados encontrados possam servir de apoio para pesquisas que busquem identificar mais a fundo os fatores que afetam a oferta de cursos no brasil

Existe uma distinção clara entre as áreas dos cursos de graduação

Para ajudar a refinar a busca por esses textos nas bases de dados (como SciELO ou Periódicos Capes), você pode utilizar combinações de palavras-chave como: "Censo da Educação Superior" AND "ensino público e privado" AND "oferta de cursos"

A Contextualização: Apresentar o panorama da educação superior no Brasil, o papel do Censo e o crescimento das redes pública e privada.

A Lacuna (A Justificativa): Mostrar que, embora a expansão seja conhecida, há dinâmicas específicas sobre a escolha das áreas de conhecimento que precisam de maior investigação.

Educação superior em turismo no Brasil: Análise da oferta de
cursos superiores no Nordeste brasileiro pelos institutos
federais

O setor privado de
ensino superior no
Brasil: continuidades
e transformações

A Pergunta: Inserir o questionamento como o elo que unirá o que já se sabe ao que o seu estudo pretende descobrir.

\section{Base de Dados}

Para a tentar responder a pergunta proposta, utilizou-se a base de dados do Censo da Educação Superior de 2024, pesquisa estatística anual coordenada pelo Instituto Nacional de Estudos e Pesquisas Educacionais Anísio Teixeira (INEP) em articulação com as Instituições de Educação Superior (IES) brasileiras. Os microdados foram coletados a partir do cadastro do Sistema e-MEC, abrangendo informações sobre instituições, cursos, infraestrutura, vagas, matrículas, ingressantes, concluintes e docentes, sendo amplamente reconhecido como o instrumento mais completo de caracterização do sistema de educação superior no país.

Para este estudo, foram selecionados os microdados do Cadastro de Cursos disponibilizados publicamente pelo INEP. O arquivo original contém um total de 720.349 observações (registros de cursos), descritas por um conjunto de 223 variáveis. A população alvo da pesquisa foram todos os cursos de graduação registrados no Brasil e que declararam informações ao censo no ano de referência. Ressalta-se os resultados das análises refletem a base após a aplicação dos filtros de seleção descritos a seguir.

\subsection{Critérios de Seleção e Variáveis Analisadas}

O processo de amostragem e delimitação da base foi conduzido em duas etapas principais:

\begin{enumerate}
    \item Filtro por nível acadêmico: foram retidos exclusivamente os registros referentes a cursos de graduação (bacharelado, licenciatura e tecnólogos), excluindo-se automaticamente todos os cursos de pós-graduação (lato e stricto sensu). Este procedimento reduziu o escopo da análise ao objeto central da pesquisa.
    \item Seleção de variáveis: das 223 variáveis disponíveis, optou-se por um conjunto reduzido e focado nas dimensões de interesse da pesquisa. As variáveis selecionadas, com suas respectivas definições operacionais, são apresentadas na Tabela 1.
\end{enumerate}

\begin{center}
\begin{tabularx}{\linewidth}{ X  X }
  \hline
  \thead{Variável} & \thead{Descrição / Categoria analisada} \\
  \hline
Rede de Ensino & Classificação binária da IES como pública ou privada. \\
  \hline
Categoria Administrativa & Subdivisão detalhada da rede (ex.: Federal, Estadual, Municipal, Privada com fins lucrativos, Privada sem fins lucrativos). \\
  \hline
Área Geral & Nome da área geral do curso, conforme adaptação brasileira da Classificação Internacional Normalizada da Educação (ISCED-F) da UNESCO. \\
  \hline
Área Específica & Nome da área específica do curso (segundo nível hierárquico da classificação CINE). \\
  \hline
Área Detalhada & Nome da área detalhada do curso (terceiro nível hierárquico da classificação CINE). \\
  \hline
Nível Acadêmico & Tipo de curso (utilizado apenas para filtrar cursos de graduação). \\
  \hline
Código do Curso & Código único de identificação do curso no Cadastro Nacional de Cursos e Instituições da Educação Superior (Cadastro e-MEC). \\
  \hline
Nome do Curso & Denominação oficial do curso. \\
  \hline
Grau Acadêmico & Tipo de grau conferido ao concluinte (Bacharelado, Licenciatura ou Tecnológico). \\
  \hline
Modalidade de Ensino & Forma de oferta do curso (Presencial ou a Distância --- EAD). \\
  \hline
Região & Região geográfica do local de oferta do curso (Norte, Nordeste, Centro-Oeste, Sudeste, Sul). \\
  \hline
\end{tabularx}
\end{center}
Fonte: Elaboração própria com base na documentação do Censo da Educação Superior 2024 (INEP, 2025).


\section{Análise exploratória}


\subsection{Distribuição das áreas gerais por rede}

Do total de 720.349 ofertas de cursos registradas na base, 700.198 (97,20\%) pertencem à rede privada e 20.151 (2,80\%) à rede pública.

A Tabela~\ref{tab:area-rede} apresenta a distribuição das áreas gerais CINE dentro de cada rede, ordenada pela diferença em pontos percentuais (p.p.) entre pública e privada.

\begin{table}[h]
\centering
\small
\begin{tabular}{lrrr}
\toprule
Área CINE Geral & Privada (\%) & Pública (\%) & Dif. (Pub $-$ Priv) \\
\midrule
Educação & 19,63 & 43,65 & $+$24,02 \\
Engenharia, produção e construção & 8,94 & 12,83 & $+$3,89 \\
Agricultura, silvicultura, pesca e vet. & 1,09 & 3,46 & $+$2,37 \\
Ciências naturais, matemática e estat. & 1,33 & 3,37 & $+$2,04 \\
Programas básicos & 0,25 & 1,32 & $+$1,07 \\
Ciências sociais, comunicação e inf. & 4,46 & 3,28 & $-$1,18 \\
Computação e TIC & 10,21 & 8,20 & $-$2,01 \\
Artes e humanidades & 5,90 & 3,21 & $-$2,69 \\
Serviços & 7,01 & 1,29 & $-$5,72 \\
Saúde e bem-estar & 11,56 & 4,78 & $-$6,78 \\
Negócios, administração e direito & 29,62 & 14,61 & $-$15,01 \\
\bottomrule
\end{tabular}
\caption{Distribuição percentual das áreas gerais CINE dentro de cada rede de ensino. Diferenças positivas indicam áreas proporcionalmente mais presentes na rede pública; negativas, na privada.}
\label{tab:area-rede}
\end{table}

A diferença entre os cursos mais populares das redes: o curso típico da rede privada está em \textit{Negócios, administração e direito} (29,6\%), enquanto o da rede pública está em \textit{Educação} (43,6\%). A rede pública é, no conjunto, mais concentrada (Educação responde por quase metade do total), enquanto a privada distribui-se entre Negócios, Educação, Saúde, TIC e Serviços.

Em termos de perfil:
\begin{itemize}
    \item \textbf{Áreas com maior peso na rede pública:} Educação ($+$24,0~p.p.), Engenharias ($+$3,9), Agricultura/Veterinária ($+$2,4), Ciências naturais e matemática ($+$2,0), Programas básicos ($+$1,1). Configura-se um perfil voltado à formação de professores, ciências básicas e áreas tradicionalmente vinculadas à pesquisa.
    \item \textbf{Áreas com maior peso na rede privada:} Negócios/Direito ($-$15,0~p.p.), Saúde ($-$6,8), Serviços ($-$5,7), Artes/Humanidades ($-$2,7), Computação/TIC ($-$2,0). Perfil compatível com graduações de forte demanda no mercado de trabalho.
    \item \textbf{Área mais neutra:} Computação e TIC, com participação próxima nas duas redes (10,2\% vs.\ 8,2\%).
\end{itemize}

A Figura~\ref{fig:area-geral-bi} apresenta visualmente o contraste entre as distribuições por rede. %e a Figura~\ref{fig:area-geral-heat} traz a mesma informação em formato de heatmap, evidenciando a magnitude das assimetrias.

\begin{figure}[h]
\centering
\includegraphics[width=0.85\textwidth]{descritiva_plots/bi_NO_CINE_AREA_GERAL_por_rede.png}
\caption{Distribuição percentual das áreas gerais CINE dentro de cada rede de ensino (\% por coluna). Fonte: elaboração própria a partir do Censo da Educação Superior 2024 (INEP).}
\label{fig:area-geral-bi}
\end{figure}

% \begin{figure}[h]
% \centering
% \includegraphics[width=0.7\textwidth]{descritiva_plots/heatmap_NO_CINE_AREA_GERAL_por_rede.png}
% \caption{Heatmap de proporções (\% dentro de cada rede) para as áreas gerais CINE. Tons mais escuros indicam maior concentração da rede correspondente naquela área.}
% \label{fig:area-geral-heat}
% \end{figure}

% \subsection{Áreas específicas: confirmação e refinamento}

% Ao descer para o segundo nível CINE (\texttt{NO\_CINE\_AREA\_ESPECIFICA}, 39 categorias), o padrão da Tabela~\ref{tab:area-rede} se confirma e ganha nuances. A área específica \textit{Educação} concentra 43,4\% da rede pública contra 19,3\% da privada, replicando a assimetria observada no nível geral. Em \textit{Negócios e administração}, a relação se inverte: 27,9\% da privada contra 13,4\% da pública.

% Dois pontos merecem destaque: (i) \textit{Serviços de segurança} aparece como área quase exclusivamente privada (3,1\% vs.\ 0,3\%); (ii) áreas como \textit{Silvicultura}, \textit{Pesca}, \textit{Meio ambiente} e \textit{Matemática e estatística}, embora pequenas em volume, são quase exclusivamente públicas. Além disso, no nível específico observa-se que \textit{Computação e TIC} aparece mais concentrada na privada (9,0\% vs.\ 4,2\%), enquanto a pública distribui parte da oferta em \textit{Programas interdisciplinares abrangendo TIC} (4,0\% vs.\ 1,2\%) --- indicando que a comparação apenas no nível geral pode subestimar diferenças mascaradas por nomenclaturas interdisciplinares.

\subsection{Variáveis complementares}

A análise das áreas não pode ser dissociada de outras dimensões em que as redes diferem fortemente.

\paragraph{Grau acadêmico.} A rede privada é majoritariamente composta por tecnólogos (48,0\%), enquanto a pública divide-se quase igualmente entre bacharelado (42,8\%) e licenciatura (43,6\%). Como licenciatura concentra-se em \textit{Educação} e cursos tecnólogos concentram-se em \textit{Negócios} e \textit{TIC}, parte relevante da diferença observada em áreas pode ser absorvida pelo grau. A Figura~\ref{fig:grau-bi} ilustra essa assimetria.

\begin{figure}[h]
\centering
\includegraphics[width=0.8\textwidth]{descritiva_plots/bi_TP_GRAU_ACADEMICO_por_rede.png}
\caption{Distribuição do grau acadêmico (Bacharelado, Licenciatura, Tecnológico) dentro de cada rede de ensino.}
\label{fig:grau-bi}
\end{figure}

\paragraph{Modalidade de ensino.} A discrepância é dramática: 96,6\% da oferta privada está em EAD, contra 45,2\% na pública (Figura~\ref{fig:modalidade-bi}). Como EAD se concentra em certas áreas (TIC, Negócios, Educação a distância), a modalidade é um segundo eixo a ser controlado.

\begin{figure}[h]
\centering
\includegraphics[width=0.6\textwidth]{descritiva_plots/bi_TP_MODALIDADE_ENSINO_por_rede.png}
\caption{Distribuição da modalidade de ensino (Presencial vs.\ EAD) dentro de cada rede.}
\label{fig:modalidade-bi}
\end{figure}


\subsection{Geografia}

A distribuição regional (Figura~\ref{fig:regiao-bi}) mostra que o Sudeste é o eixo principal em ambas as redes (cerca de 40\% das ofertas). As diferenças regionais marcantes são: \textit{Sul mais privado} (22,1\% vs.\ 13,7\%) e \textit{Nordeste mais público} (25,6\% vs.\ 19,6\%).

\begin{figure}[h]
\centering
\includegraphics[width=0.7\textwidth]{descritiva_plots/bi_NO_REGIAO_por_rede.png}
\caption{Distribuição regional da oferta de cursos dentro de cada rede de ensino.}
\label{fig:regiao-bi}
\end{figure}

\subsection{Síntese exploratória e hipóteses para etapas seguintes}

A partir do exposto, e ainda em caráter puramente descritivo, observa-se:

\begin{enumerate}
    \item \textbf{Há sinais robustos de distinção entre as áreas das duas redes.} A moda muda (Negócios na privada, Educação na pública), e diferenças passam de 20~p.p.\ em pelo menos uma área.
    \item \textbf{A distinção segue um perfil reconhecível:} privadas concentram cursos com forte vínculo ao mercado de trabalho imediato (Negócios, Saúde, Serviços, TIC); públicas concentram-se em formação de professores (Educação/Licenciaturas) e áreas-base do conhecimento (Ciências naturais, Engenharias, Agricultura).
    \item \textbf{Essa distinção está entrelaçada com outras dimensões institucionais} --- sobretudo grau acadêmico e modalidade de ensino --- que não são "ruído", mas parte do mesmo fenômeno (modelo de negócio das privadas vs.\ missão institucional das públicas).
\end{enumerate}

Estes padrões motivam um conjunto de hipóteses a serem investigadas nas seções de Metodologia e Resultados:

\begin{itemize}
    \item \textbf{H1.} A distribuição de áreas gerais CINE não é independente da rede de ensino.
    \item \textbf{H2.} A diferença concentra-se principalmente em duas áreas: \textit{Educação} (mais pública) e \textit{Negócios, administração e direito} (mais privada).
    \item \textbf{H3.} Após considerar o grau acadêmico, parte da diferença em áreas se atenua, mas a assimetria em \textit{Educação} persiste pela ligação com Licenciatura.
    \item \textbf{H4.} Após considerar a modalidade de ensino, a assimetria entre redes permanece, dado que a oferta privada em EAD continua concentrada em Negócios e Educação a distância.
    \item \textbf{H5.} Dentro da rede privada, instituições comunitárias/confessionais apresentam distribuição de áreas mais próxima da rede pública do que privadas com fins lucrativos.
    \item \textbf{H6.} O perfil de áreas por rede é estável entre as regiões do país.
\end{itemize}


\section{Introdução}

Apresente o contexto do problema estudado. Explique por que o tema é relevante e qual pergunta de pesquisa o trabalho pretende responder.

Exemplo de perguntas:

\begin{itemize}
\item Existe associação entre escolaridade e tipo de ocupação?
\item O sexo influencia a probabilidade de determinado evento?
\end{itemize}

Existe trabalho relacionado na área? Faça uma pequena revisão bibliográfica

\bigskip
Finalize a introdução com o objetivo do estudo.

\section{Dados}

Descreva a base de dados utilizada:

\begin{itemize}
\item origem da base
\item número de observações
\item variáveis analisadas
\item possíveis limitações
\end{itemize}

Se necessário, apresente tabelas descritivas e gráficos.

\begin{table}[h]
\centering
\begin{tabular}{ccc}
\toprule
Categoria & Frequência & Proporção \\
\midrule
A & 100 & 0.40 \\
B & 80 & 0.32 \\
C & 70 & 0.28 \\
\bottomrule
\end{tabular}
\caption{Exemplo de tabela descritiva}
\end{table}

\section{Metodologia}

Descreva os métodos estatísticos utilizados.

Exemplos:

\begin{itemize}
\item Teste qui-quadrado de independência
\item Teste exato de Fisher
\item Razão de chances (odds ratio)
\item Regressão logística
\item Modelos log-lineares
\end{itemize}

Apresente os modelos quando necessário.

Exemplo de regressão logística:

\begin{equation}
\log\left(\frac{p}{1-p}\right) = \beta_0 + \beta_1 x_1 + \beta_2 x_2
\end{equation}

Explique brevemente como os parâmetros são interpretados.

\section{Resultados}

Apresente os resultados obtidos.

Inclua:

\begin{itemize}
\item tabelas
\item gráficos
\item estimativas dos modelos
\item interpretação estatística
\end{itemize}

Exemplo de figura:

\begin{figure}[h]
\centering
\includegraphics[width=0.7\textwidth]{figura_exemplo}
\caption{Exemplo de gráfico}
\end{figure}

\section{Discussão}

Interprete os resultados no contexto do problema.

Discuta:

\begin{itemize}
\item implicações dos resultados
\item possíveis limitações da análise
\item sugestões para estudos futuros
\end{itemize}

\section{Conclusão}

Resuma os principais achados do trabalho e retome a pergunta de pesquisa.

\section*{Referências}

Inclua referências utilizadas no trabalho.

Exemplo:

Agresti, A. (2013). \textit{Categorical Data Analysis}. Wiley.

\section*{Apêndice (opcional)}

Trechos de código em R podem ser incluídos aqui.

\end{document}